\documentclass[10pt,a4paper]{amsart}
\usepackage[margin=19mm]{geometry}
\usepackage{amsmath,amssymb,mathtools}
\usepackage[scr=rsfs,cal=euler]{mathalfa}
\usepackage{xfrac}
\usepackage[unicode,hidelinks,bookmarks=false]{hyperref}
\newcommand{\ie}{i.e.\ }
\newcommand{\cf}{cf.\ }
\newcommand{\resp}{respectively\ }
\newcommand{\Subsection}{\subsection}
\providecommand{\nobreakdash}{\nobreak\hbox{-}}
\setlength{\emergencystretch}{2em}
\multlinegap0pt
\usepackage{xcolor}
\usepackage{float}
\usepackage{algorithm}\usepackage{algpseudocode}
\usepackage[shortlabels]{enumitem}
\newtheorem{lemma}{Lemma}
\newtheorem{thm}{Theorem}
\title{\LARGE \bf Soft projections for robust data-driven control}
\author{Andr\'as Sasfi, Jaap Eising, Florian D\"orfler}
\date{Source: arXiv:2604.00900v1 (CC BY 4.0)}
\begin{document}
\maketitle

\noindent\emph{Source and licence.} \LARGE \bf Soft projections for robust data-driven control. Version arXiv:2604.00900v1 (CC BY 4.0), distributed by arXiv under the Creative Commons Attribution 4.0 International licence, http://creativecommons.org/licenses/by/4.0/, as recorded in the arXiv licence metadata for this exact version and shown on the abstract page https://arxiv.org/abs/2604.00900v1. Full licence notice: \texttt{licenses/arxiv-2604.00900-CC-BY-4.0.txt}.\par\smallskip

\noindent\textit{Editorial scope.} Counted unit: the article's thm thm:generalizing\_DeePC (source0.tex lines 364--368). The article's complete original proof of it is delivered with it (source0.tex lines 370--393, one \texttt{proof} environment of 24 lines). No mathematics is written, repaired or re-derived: the statement and the proof are the article's own lines, in the article's own notation, and no retained line is substituted or re-flowed.  Retained uncounted as the source-local context the proof needs: lines 188--196; lines 206--216; lines 358--362.  Nothing was omitted from any retained span, so the package carries no omission record.  No retained line contains a citation, so the wrapper carries no bibliography and no bibliography entry.  No retained line contains a figure environment, an image-inclusion command or drawing code, so no image is delivered and the package declares no asset dependency.  Editorial formatting only, in addition: an AMS \texttt{amsart} preamble that restates the article's own declarations and macros used by the retained lines, namely the environments \texttt{lemma}, \texttt{thm} on separate counters, together with the standard packages those lines need.  The retained environments print in this document as Lemma 1, Theorem 1 in order of appearance, whereas the article prints the same items with the numbering of its own full document.  The complete native source of the article as distributed in the arXiv e-print for this exact version is bundled unchanged as \texttt{sources/197619/source0.tex}.
\begin{align} \label{eq:DeePC}
    \begin{split}
    \min_{w \in \mathcal{C},g\in\mathbb{R}^D} & \quad J(w_\mathrm{f}) + \lambda_\sigma \|\sigma\|_2^2 + \lambda_g\|g\|_2^2 \\
    \mathrm{s.t.} & \quad Hg = \begin{bmatrix}
        \hat{w}_\mathrm{ini} + \sigma \\ 
        w_\mathrm{f}
    \end{bmatrix}.
    \end{split}
\end{align}

\begin{lemma} \label{lemma:DeePC}
    If $\mathcal{C} = \mathbb{R}^{qL}$ and $\lambda_g>0$, the solution to~\eqref{eq:DeePC} is
    \begin{align*}
        w^\star_\mathrm{DeePC} & =  H(H^\top W H + \lambda_g I)^{-1} H^\top W \begin{bmatrix}
    \hat{w}_\mathrm{ini} \\ w_\mathrm{ref}
\end{bmatrix}, \\
    g^\star & = (H^\top W H + \lambda_g I)^{-1}H^\top W \begin{bmatrix}
            \hat{w}_\mathrm{ini} \\ w_\mathrm{ref}
        \end{bmatrix}.
    \end{align*}
\end{lemma}

\begin{align} \label{eq:indicator_approx}
     \min_{w\in\mathcal{C}} \left \|w - \begin{bmatrix}
    \hat{w}_\mathrm{ini} \\ w_\mathrm{ref}
    \end{bmatrix}\right \|^2_{W} + \frac{\hat{\alpha}}{\delta} w^\top(I - \tilde{P}_H)w. 
\end{align}

\begin{thm} \label{thm:generalizing_DeePC}
   Let $w^\star_\mathrm{DeePC}$ and $w^\star_\eqref{eq:indicator_approx}$ be the solutions to~\eqref{eq:DeePC} and \eqref{eq:indicator_approx}, respectively, with $\mathcal{C}=\mathbb{R}^{qL}$.
   Assume that $H$ is full row rank and $\lambda_g,\delta>0$.
   If $\hat{\alpha} = \lambda_g$, then $w^\star_\eqref{eq:indicator_approx} \to w^\star_\mathrm{DeePC}$ as $\delta \to 0$. 
\end{thm}

\begin{proof}
    We first reformulate the solution $w^\star_\mathrm{DeePC}$ from Lemma~\ref{lemma:DeePC}.
    Using the identity $(I +AB)^{-1}A = A(I+BA)^{-1}$ with $A = H^\top$ and $B=\lambda_g^{-1}WH$ leads to
    \begin{align*}
        H(H^\top W H + & \lambda_g I)^{-1} H^\top  = \lambda_g^{-1} H(\lambda_g^{-1} H^\top W H + I)^{-1}H^\top \\
        &= \lambda_g^{-1} H H^\top(\lambda_g^{-1}WHH^\top+I)^{-1} \\
        & = \left((\lambda_g^{-1}WHH^\top +I)\cdot\lambda_g(HH^\top)^{-1}\right)^{-1} \\
        & = \left(W + \lambda_g(HH^\top)^{-1}\right)^{-1}.
    \end{align*}
    Thus, $w^\star_\mathrm{DeePC} = \left(W + \lambda_g(HH^\top)^{-1}\right)^{-1}W\begin{bmatrix}
        \hat{w}_\mathrm{ini}^\top & w_\mathrm{ref}^\top
    \end{bmatrix}^\top.$
    Furthermore,~\eqref{eq:indicator_approx} without constraints is a regularized least-squares problem, whose solution is
    \begin{align*}
    w^\star_\eqref{eq:indicator_approx} &= \left(W + \hat{\alpha}/\delta \cdot(I-\tilde{P}_H)\right)^{-1}W\begin{bmatrix}
        \hat{w}_\mathrm{ini}^\top & w_\mathrm{ref}^\top
    \end{bmatrix}^\top \\
    & = \left(W + \hat{\alpha}(HH^\top+\delta I)^{-1}\right)^{-1}W\begin{bmatrix}
        \hat{w}_\mathrm{ini}^\top & w_\mathrm{ref}^\top
    \end{bmatrix}^\top.
    \end{align*}
    Note that the inverses exist for any $\delta>0$, since $W$ is positive definite, $\hat{\alpha}>0$, and $H$ is full row rank.
    Therefore, if $\hat{\alpha}=\lambda_g$, then $w^\star_\eqref{eq:indicator_approx} \to w^\star_\mathrm{DeePC}$ as $\delta\to0$.
\end{proof}

\end{document}
